\documentclass[12pt,a4paper]{article}

\usepackage[margin=2.5cm]{geometry}

\usepackage{fontspec}
\setmainfont{lmroman10-regular.otf}[
  BoldFont=lmroman10-bold.otf,
  ItalicFont=lmroman10-italic.otf,
  BoldItalicFont=lmroman10-bolditalic.otf]
\setmonofont{lmmono10-regular.otf}[
  BoldFont=lmmonolt10-bold.otf,
  ItalicFont=lmmono10-italic.otf]

\usepackage{polyglossia}
\setdefaultlanguage{vietnamese}
\setotherlanguage{english}

\usepackage{xcolor}
\usepackage{colortbl}
\usepackage{array}
\usepackage{longtable}
\usepackage{enumitem}
\usepackage{titlesec}
\usepackage{fancyhdr}
\setlength{\headheight}{15pt}
\usepackage[most]{tcolorbox}
\usepackage{fancyvrb}
\usepackage{hyperref}
\usepackage{bookmark}

\hypersetup{
  colorlinks=true,
  linkcolor=blue!55!black,
  urlcolor=blue!55!black,
  pdftitle={Chuong 5: Nguoi dung va phan quyen},
  pdfauthor={Tai lieu Linux/DevOps}
}

% ---------- Mau ----------
\definecolor{labcolor}{HTML}{1B6CA8}
\definecolor{labbg}{HTML}{EAF3FA}
\definecolor{seniorcolor}{HTML}{6B4FA0}
\definecolor{seniorbg}{HTML}{F1ECF9}
\definecolor{checkcolor}{HTML}{1F8A52}
\definecolor{checkbg}{HTML}{E8F6EE}
\definecolor{notecolor}{HTML}{C77C00}
\definecolor{notebg}{HTML}{FBF2E0}
\definecolor{warncolor}{HTML}{B3261E}
\definecolor{warnbg}{HTML}{FBE9E7}
\definecolor{headrule}{HTML}{888888}
\definecolor{tablehead}{HTML}{D9E6F2}

% ---------- Tieu de muc ----------
\titleformat{\section}
  {\Large\bfseries\color{labcolor!80!black}}{\thesection}{0.6em}{}
\titleformat{\subsection}
  {\large\bfseries\color{seniorcolor!75!black}}{\thesubsection}{0.6em}{}

% ---------- Header/Footer ----------
\pagestyle{fancy}
\fancyhf{}
\renewcommand{\headrulewidth}{0.6pt}
\renewcommand{\footrulewidth}{0pt}
\fancyhead[L]{\small\itshape Chương 5 — Người dùng \& phân quyền}
\fancyhead[R]{\small\thepage}

% ---------- Verbatim style ----------
\fvset{fontsize=\small,frame=single,framesep=4pt,rulecolor=\color{headrule},xleftmargin=4pt,xrightmargin=4pt}

% ---------- Cac hop tcolorbox ----------
\newtcolorbox{labbox}[1][Thực hành ngay]{
  enhanced, breakable,
  colback=labbg, colframe=labcolor,
  coltitle=white, fonttitle=\bfseries,
  title={#1}, left=3mm, right=3mm, top=2mm, bottom=2mm,
  boxrule=0.8pt, arc=2mm}

\newtcolorbox{seniorbox}[1][Góc Senior]{
  enhanced, breakable,
  colback=seniorbg, colframe=seniorcolor,
  coltitle=white, fonttitle=\bfseries,
  title={#1}, left=3mm, right=3mm, top=2mm, bottom=2mm,
  boxrule=0.8pt, arc=2mm}

\newtcolorbox{checkbox}[1][Tự kiểm tra]{
  enhanced, breakable,
  colback=checkbg, colframe=checkcolor,
  coltitle=white, fonttitle=\bfseries,
  title={#1}, left=3mm, right=3mm, top=2mm, bottom=2mm,
  boxrule=0.8pt, arc=2mm}

\newtcolorbox{notebox}[1][Ghi chú]{
  enhanced, breakable,
  colback=notebg, colframe=notecolor,
  coltitle=white, fonttitle=\bfseries,
  title={#1}, left=3mm, right=3mm, top=2mm, bottom=2mm,
  boxrule=0.8pt, arc=2mm}

\newtcolorbox{warnbox}[1][Cảnh báo]{
  enhanced, breakable,
  colback=warnbg, colframe=warncolor,
  coltitle=white, fonttitle=\bfseries,
  title={#1}, left=3mm, right=3mm, top=2mm, bottom=2mm,
  boxrule=1pt, arc=2mm}

\setlist[itemize]{leftmargin=1.4em,topsep=2pt,itemsep=2pt}
\setlist[enumerate]{leftmargin=1.6em,topsep=2pt,itemsep=2pt}

\newcommand{\cmd}[1]{\texttt{#1}}

\begin{document}

% ================= TRANG TIEU DE =================
\begin{titlepage}
\centering
\vspace*{3cm}
{\Huge\bfseries\color{labcolor!80!black} Chương 5\par}
\vspace{1cm}
{\LARGE\bfseries Người dùng \& phân quyền\par}
\vspace{0.4cm}
{\Large ai được làm gì (\texttt{chmod}, \texttt{sudo})\par}
\vspace{2cm}
\begin{tcolorbox}[colback=labbg,colframe=labcolor,width=0.85\textwidth,arc=3mm,boxrule=0.8pt]
\itshape
Nối tiếp Chương 1--4: bạn đã biết \cmd{cd}, \cmd{ls}, \cmd{file}, chuyển hướng
(\cmd{>}), đường ống (\cmd{|}), và đã từng nhìn thấy cột quyền lạ lạ kiểu
\cmd{-rw-r-{}-r-{}-} khi gõ \cmd{ls -l}. Chương này sẽ giải mã đúng cái cột đó,
và dạy bạn cách thay đổi ``ai được làm gì'' với file của mình.
\end{tcolorbox}
\vfill
{\large Tài liệu học Linux / DevOps cho người mới\par}
\vspace{0.5cm}
{\small Thực hành trên macOS (shell \texttt{zsh})\par}
\end{titlepage}

\tableofcontents
\newpage

% ================= 1. MUC TIEU =================
\section{Mục tiêu học tập}

Học xong chương này, bạn sẽ:

\begin{enumerate}
\item Hiểu vì sao Linux/Unix (và macOS) là hệ \textbf{đa người dùng}, và vì sao
  điều đó bắt buộc phải có \textbf{phân quyền}.
\item Biết mình \textbf{là ai} trên máy: dùng \cmd{whoami}, \cmd{id}, \cmd{groups}.
\item \textbf{Đọc trôi chảy} chuỗi 10 ký tự quyền trong \cmd{ls -l} (ví dụ
  \cmd{-rwxr-xr-{}-}), phân biệt quyền của \textbf{chủ sở hữu (owner)} /
  \textbf{nhóm (group)} / \textbf{người khác (others)}.
\item Hiểu ý nghĩa \textbf{r (đọc)}, \textbf{w (ghi)}, \textbf{x (thực thi)} --- và
  đặc biệt vì sao \cmd{x} trên \textbf{file} khác \cmd{x} trên \textbf{thư mục}.
\item Dùng \cmd{chmod} thành thạo theo \textbf{2 cách}: ký hiệu (\cmd{u+x}) và số
  (\cmd{755}, \cmd{644}, \cmd{600}).
\item Biết sơ về \cmd{chown}/\cmd{chgrp} (đổi chủ/nhóm) và sự khác biệt trên macOS.
\item Hiểu \cmd{sudo} là gì, khi nào cần, vì sao \textbf{nguy hiểm nếu lạm dụng},
  và khác \cmd{su} ra sao.
\item Nắm tư duy \textbf{least privilege (đặc quyền tối thiểu)} --- nền tảng bảo
  mật của mọi DevOps/Senior.
\end{enumerate}

\begin{notebox}[Trấn an]
Toàn bộ phần thực hành ở chương này \textbf{không dùng \cmd{sudo}}, \textbf{không
đụng file hệ thống}. Bạn chỉ nghịch trong sân tập \cmd{\textasciitilde/linux-lab}.
An toàn tuyệt đối, có lỡ gõ sai cũng không hỏng máy.
\end{notebox}

% ================= 2. VI SAO CAN PHAN QUYEN =================
\section{Vì sao cần phân quyền? (câu chuyện đa người dùng)}

Hãy tưởng tượng một \textbf{văn phòng dùng chung}: nhiều người ngồi chung, mỗi
người có một ngăn tủ khóa riêng. Bạn không muốn đồng nghiệp mở tủ bạn lấy đồ, và
sếp (người có chìa khóa tổng) thì mở được mọi tủ khi cần.

Unix/Linux sinh ra trong môi trường y hệt: \textbf{một máy chủ, nhiều người dùng}
đăng nhập cùng lúc, chạy nhiều tiến trình (chương trình đang chạy) song song. Để
mọi người không giẫm chân/đọc trộm/xóa nhầm của nhau, hệ điều hành gán cho
\textbf{mỗi file và thư mục} một bộ quy tắc: \textit{ai được đọc, ai được sửa, ai
được chạy}.

macOS kế thừa nguyên vẹn triết lý này (vì macOS có lõi Unix BSD). Dù máy bạn chỉ có
một mình bạn dùng, hệ điều hành vẫn âm thầm chạy hàng chục \textbf{tài khoản hệ
thống} (cho tiến trình nền) --- nên phân quyền vẫn hoạt động đầy đủ.

\textbf{Ba ``vai'' cần nhớ} với mỗi file:

\begin{center}
\begin{tabular}{|>{\raggedright\arraybackslash}p{3.2cm}|>{\raggedright\arraybackslash}p{3.8cm}|>{\raggedright\arraybackslash}p{5.5cm}|}
\hline
\rowcolor{tablehead}\textbf{Vai} & \textbf{Tiếng Anh} & \textbf{Là ai} \\
\hline
Chủ sở hữu & \textbf{owner} (hay \textit{user}) & Người tạo/sở hữu file \\
\hline
Nhóm & \textbf{group} & Một nhóm tài khoản; file thuộc về 1 nhóm \\
\hline
Người khác & \textbf{others} & Tất cả những ai còn lại \\
\hline
\end{tabular}
\end{center}

Và trên đỉnh kim tự tháp là \textbf{root} (còn gọi \textbf{superuser} --- siêu người
dùng): tài khoản có \textbf{quyền tối cao}, làm được mọi thứ, bỏ qua mọi rào quyền.
(Ta sẽ gặp root khi nói về \cmd{sudo} ở mục 9.)

\begin{notebox}[Định nghĩa nhanh]
\begin{itemize}
\item \textbf{uid} = User ID = số định danh người dùng (vd root luôn là uid \cmd{0}).
\item \textbf{gid} = Group ID = số định danh nhóm.
\item \textbf{tiến trình (process)} = một chương trình đang chạy.
\end{itemize}
\end{notebox}

\begin{checkbox}[Tự kiểm tra 5.1]
\begin{enumerate}
\item Ba ``vai'' quyền của một file là gì?
\item root có gì đặc biệt?
\item Vì sao máy cá nhân (chỉ mình bạn) vẫn cần phân quyền?
\end{enumerate}
\end{checkbox}

% ================= 3. LAB A =================
\section{Thực hành ngay (lab A): Tôi là ai trên máy này?}

Mở \textbf{Terminal} (ứng dụng Terminal trên macOS, shell \cmd{zsh}). Gõ từng dòng,
nhấn Enter sau mỗi dòng.

\begin{labbox}[Lab A --- Nhận diện danh tính]
\begin{Verbatim}
whoami
\end{Verbatim}
\textbf{Kết quả mong đợi:} in ra \textbf{tên đăng nhập} của bạn, ví dụ:
\begin{Verbatim}
vinh
\end{Verbatim}

\begin{Verbatim}
id
\end{Verbatim}
\textbf{Kết quả mong đợi:} một dòng dài chứa uid, gid và danh sách nhóm, ví dụ:
\begin{Verbatim}
uid=501(vinh) gid=20(staff) groups=20(staff),12(everyone),61(localaccounts),...
\end{Verbatim}
\begin{itemize}
\item \cmd{uid=501(vinh)}: bạn là người dùng số 501, tên \cmd{vinh}. (Trên macOS,
  tài khoản người đầu tiên thường bắt đầu từ \textbf{501}.)
\item \cmd{gid=20(staff)}: nhóm chính (primary group) \textbf{của phiên làm việc
  hiện tại} là \textbf{staff} (số 20). Đây là điểm \textbf{đặc trưng của macOS} ---
  trên Linux, nhóm chính thường trùng tên người dùng.
\item \cmd{groups=...}: tất cả các nhóm bạn thuộc về.
\end{itemize}

\begin{Verbatim}
groups
\end{Verbatim}
\textbf{Kết quả mong đợi:} liệt kê tên các nhóm, ví dụ:
\begin{Verbatim}
staff everyone localaccounts ...
\end{Verbatim}
\end{labbox}

\begin{notebox}[Máy bạn có thể khác]
Tên đăng nhập, số uid, và danh sách nhóm sẽ khác --- đó là bình thường. Quan trọng
là bạn thấy được \textbf{3 lệnh này trả lời câu ``tôi là ai''} như thế nào. Nếu bạn
là tài khoản đầu tiên trên Mac, gần như chắc chắn nhóm chính là \cmd{staff}.
\end{notebox}

\begin{notebox}[Góc thuật ngữ]
\cmd{whoami} = ``who am I'' (tôi là ai). \cmd{id} = identity (danh tính). Dễ nhớ.
\end{notebox}

\begin{checkbox}[Tự kiểm tra 5.2]
\begin{enumerate}
\item Lệnh nào cho biết \textbf{tên đăng nhập} của bạn?
\item Trong \cmd{id}, \cmd{uid} và \cmd{gid} khác nhau ở chỗ nào?
\item Nhóm chính mặc định trên macOS thường tên là gì?
\end{enumerate}
\end{checkbox}

% ================= 4. DOC COT QUYEN =================
\section{Đọc cột quyền trong \texttt{ls -l} (giải mã 10 ký tự)}

Đây là kỹ năng \textbf{quan trọng nhất} của chương. Khi gõ \cmd{ls -l}, mỗi dòng bắt
đầu bằng một chuỗi \textbf{10 ký tự}, ví dụ:

\begin{Verbatim}
-rwxr-xr--
\end{Verbatim}

Ta chia nó thành \textbf{1 + 3 + 3 + 3}:

\begin{Verbatim}
 -        rwx        r-x        r--
 ^         ^          ^          ^
loai    owner      group      others
file   (chu)      (nhom)    (nguoi khac)
\end{Verbatim}

\subsection{Ký tự đầu tiên: loại đối tượng}

\begin{center}
\begin{tabular}{|c|>{\raggedright\arraybackslash}p{9cm}|}
\hline
\rowcolor{tablehead}\textbf{Ký tự} & \textbf{Nghĩa} \\
\hline
\cmd{-} & file thường (regular file) \\
\hline
\cmd{d} & thư mục (\textbf{d}irectory) \\
\hline
\cmd{l} & liên kết tượng trưng (\textbf{l}ink --- như ``shortcut'') \\
\hline
\end{tabular}
\end{center}

\begin{notebox}[Còn vài loại khác ít gặp]
\cmd{c} và \cmd{b} (thiết bị --- character/block device), \cmd{p} (named pipe/FIFO),
\cmd{s} (socket). Khi mới học bạn gần như không gặp; chỉ cần biết tên để không bất
ngờ khi sau này thấy chúng ở các thư mục hệ thống.
\end{notebox}

\begin{notebox}[Đặc trưng macOS --- dấu \texttt{@} và \texttt{+}]
Trên macOS bạn có thể thấy thêm một dấu ngay \textbf{sau} 10 ký tự, ví dụ
\cmd{-rw-r-{}-r-{}-@}. Dấu \cmd{@} nghĩa là file có \textbf{thuộc tính mở rộng
(extended attributes)} --- KHÔNG phải quyền r/w/x, cứ bỏ qua khi mới học (muốn xem
chi tiết: \cmd{xattr -l tenfile}). Dấu \cmd{+} (nếu thấy) nghĩa là file có
\textbf{ACL} (danh sách quyền chi tiết hơn). Cả hai đều không ảnh hưởng tới bài học
về r/w/x.
\end{notebox}

\subsection{Ba cụm \texttt{rwx}: quyền cho từng vai}

Mỗi cụm có \textbf{3 vị trí cố định}: \cmd{r} rồi \cmd{w} rồi \cmd{x}. Nếu có quyền
thì hiện chữ, không có thì hiện dấu \cmd{-}.

\begin{center}
\begin{tabular}{|c|>{\raggedright\arraybackslash}p{6cm}|c|}
\hline
\rowcolor{tablehead}\textbf{Chữ} & \textbf{Tên} & \textbf{Giá trị số} \\
\hline
\cmd{r} & \textbf{read} --- đọc & \textbf{4} \\
\hline
\cmd{w} & \textbf{write} --- ghi/sửa & \textbf{2} \\
\hline
\cmd{x} & e\textbf{x}ecute --- thực thi & \textbf{1} \\
\hline
\cmd{-} & không có quyền đó & 0 \\
\hline
\end{tabular}
\end{center}

Ví dụ giải mã \cmd{-rwxr-xr-{}-}:

\begin{center}
\begin{tabular}{|c|c|>{\raggedright\arraybackslash}p{7cm}|}
\hline
\rowcolor{tablehead}\textbf{Cụm} & \textbf{Ký tự} & \textbf{Nghĩa} \\
\hline
Loại & \cmd{-} & file thường \\
\hline
owner & \cmd{rwx} & chủ đọc + ghi + chạy được \\
\hline
group & \cmd{r-x} & nhóm đọc + chạy được, KHÔNG sửa được \\
\hline
others & \cmd{r-{}-} & người khác CHỈ đọc được \\
\hline
\end{tabular}
\end{center}

\subsection{\texttt{x} trên FILE khác \texttt{x} trên THƯ MỤC --- đừng nhầm!}

Đây là chỗ người mới hay rối:

\begin{itemize}
\item \textbf{\cmd{x} trên một FILE} = ``file này \textbf{chạy được} như một
  chương trình/script''. Không có \cmd{x} thì dù file là script đúng, hệ thống vẫn
  \textbf{từ chối chạy}.
\item \textbf{\cmd{x} trên một THƯ MỤC} = ``được phép \textbf{đi vào} (cd) thư mục đó
  và truy cập file bên trong''. Một thư mục có \cmd{r} nhưng không có \cmd{x} thì bạn
  \textbf{liệt kê được tên file nhưng không vào được}.
\end{itemize}

\begin{notebox}[Ghi nhớ một câu]
Với thư mục, \cmd{x} nghĩa là \textit{``được bước qua cửa''}; với file, \cmd{x}
nghĩa là \textit{``được khởi động nó''}.
\end{notebox}

\begin{checkbox}[Tự kiểm tra 5.3]
\begin{enumerate}
\item Trong \cmd{-rw-r-{}-r-{}-}, \textbf{owner} có những quyền nào? \textbf{others}
  có quyền gì?
\item \cmd{r=?}, \cmd{w=?}, \cmd{x=?} (giá trị số).
\item Một thư mục có quyền \cmd{r} nhưng thiếu \cmd{x} thì bạn làm được gì và không
  làm được gì?
\end{enumerate}
\end{checkbox}

% ================= 5. LAB B =================
\section{Thực hành ngay (lab B): Quyền \texttt{x} ``sống động'' --- chạy một script}

Giờ ta sẽ tận mắt thấy quyền \cmd{x} quyết định ``chạy được hay không''. \textbf{Gõ
tay} từng dòng.

\begin{labbox}[Lab B --- Bước 1: Tạo sân tập]
\begin{Verbatim}
mkdir -p ~/linux-lab/ch5
cd ~/linux-lab/ch5
\end{Verbatim}
\begin{itemize}
\item \cmd{mkdir -p}: tạo thư mục (cờ \cmd{-p} = tạo cả thư mục cha nếu chưa có, và
  không báo lỗi nếu đã tồn tại).
\item \cmd{cd}: vào thư mục đó.
\end{itemize}
\textbf{Kết quả mong đợi:} không in gì ra cả (im lặng = thành công). Kiểm chứng bằng:
\begin{Verbatim}
pwd
\end{Verbatim}
in ra đường dẫn, ví dụ \cmd{/Users/vinh/linux-lab/ch5}.
\end{labbox}

\begin{labbox}[Lab B --- Bước 2: Tạo một script nhỏ]
\begin{Verbatim}
echo 'echo "xin chao tu script"' > chao.sh
ls -l chao.sh
\end{Verbatim}
\textbf{Kết quả mong đợi (trên macOS, chú ý dấu \texttt{@}):}
\begin{Verbatim}
-rw-r--r--@  1 vinh  staff  26 Jun 26 10:00 chao.sh
\end{Verbatim}
Đọc dòng này: cột đầu là \textbf{quyền}; sau quyền là \textbf{số liên kết}
(\cmd{1}); rồi \textbf{tên chủ sở hữu} (\cmd{vinh}); rồi \textbf{nhóm của file}
(\cmd{staff}); rồi \textbf{số byte}, ngày giờ, và tên file.

Nhìn cụm quyền: \cmd{-rw-r-{}-r-{}-}. Owner đọc+ghi, group đọc, others đọc.
\textbf{Không hề có \cmd{x} ở đâu cả} --- nghĩa là file này chưa ``chạy được''.
\end{labbox}

\begin{notebox}[Máy bạn có thể khác --- đọc kỹ để khỏi hoang mang]
\begin{itemize}
\item \textbf{Dấu \cmd{@}:} Trên macOS hiện đại, file tạo bằng \cmd{echo ... > file}
  gần như luôn nhận một thuộc tính mở rộng (vd \cmd{com.apple.provenance}), nên
  \cmd{ls -l} hiển thị dấu \cmd{@} ngay sau cụm quyền (\cmd{-rw-r-{}-r-{}-@}). Đây
  là \textbf{bình thường}, dấu \cmd{@} KHÔNG phải quyền và không ảnh hưởng gì tới
  r/w/x. Nếu máy bạn không có \cmd{@} cũng không sao.
\item \textbf{Cột group (\cmd{staff} hay \cmd{wheel}?):} Cột thứ tư có thể hiển thị
  \cmd{staff} HOẶC \cmd{wheel} --- tùy thư mục bạn đang đứng. macOS (lõi BSD) cho
  file mới \textbf{kế thừa nhóm của thư mục cha}, nên ở \cmd{\textasciitilde/linux-lab}
  thường là \cmd{staff}, còn ở vài nơi khác (vd \cmd{/tmp}) lại là \cmd{wheel}. Cả
  hai đều bình thường, không ảnh hưởng bài học.
\item \textbf{Số byte} (vd \cmd{26}), \textbf{ngày giờ}, \textbf{tên đăng nhập}
  đương nhiên khác. Phần cụm quyền \cmd{-rw-r-{}-r-{}-} thì gần như chắc chắn giống
  --- đó là quyền mặc định cho file mới tạo.
\end{itemize}
\end{notebox}

\begin{labbox}[Lab B --- Bước 3: Thử chạy khi CHƯA có \texttt{x}]
\begin{Verbatim}
./chao.sh
\end{Verbatim}
\textbf{Kết quả mong đợi:} bị từ chối, báo đại loại:
\begin{Verbatim}
zsh: permission denied: ./chao.sh
\end{Verbatim}
\textbf{Lưu ý máy bạn có thể khác:} phần tiền tố trước có thể là \cmd{zsh:} hoặc
\cmd{(eval):1:} tùy cấu hình shell --- không sao. Nếu bạn lỡ dùng \cmd{bash} thay vì
\cmd{zsh}, thông báo sẽ hơi khác: \cmd{bash: ./chao.sh: Permission denied} --- ý
nghĩa giống hệt.

Đây chính là minh họa \textbf{sống động}: file nội dung đúng, nhưng thiếu quyền
\cmd{x} nên hệ thống \textbf{không cho chạy}. (\cmd{./} nghĩa là ``file nằm ngay
trong thư mục hiện tại''.)
\end{labbox}

\begin{labbox}[Lab B --- Bước 4: Thêm quyền \texttt{x} cho owner, rồi chạy lại]
\begin{Verbatim}
chmod u+x chao.sh
ls -l chao.sh
\end{Verbatim}
\textbf{Kết quả mong đợi:} giờ có \cmd{x} ở cụm owner:
\begin{Verbatim}
-rwxr--r--@  1 vinh  staff  26 Jun 26 10:00 chao.sh
\end{Verbatim}
Chạy lại:
\begin{Verbatim}
./chao.sh
\end{Verbatim}
\textbf{Kết quả mong đợi:}
\begin{Verbatim}
xin chao tu script
\end{Verbatim}

Bạn vừa chứng kiến: \textbf{chỉ cần thêm một chữ \cmd{x}, file từ ``không chạy được''
thành ``chạy được''}. Đó là sức mạnh của phân quyền.
\end{labbox}

\begin{notebox}[Giải thích \texttt{chmod u+x}]
\cmd{chmod} = \textit{change mode} (đổi quyền). \cmd{u} = user/owner. \cmd{+x} = thêm
quyền execute. Đọc liền: ``thêm quyền chạy cho chủ sở hữu''.
\end{notebox}

\begin{notebox}[Vì sao script này chạy được]
Script chạy được là nhờ shell hiện tại (\cmd{zsh}) tự diễn giải nội dung. Ở chương
sau bạn sẽ học thêm dòng đầu \cmd{\#!/bin/sh} (gọi là \textbf{shebang}) để chỉ định
rõ trình thông dịch cho script --- khi đó file có thể chạy độc lập, không phụ thuộc
shell đang dùng.
\end{notebox}

\begin{checkbox}[Tự kiểm tra 5.4]
\begin{enumerate}
\item Vì sao lần đầu \cmd{./chao.sh} báo \cmd{permission denied}?
\item \cmd{chmod u+x chao.sh} làm gì?
\item \cmd{./} ở đầu lệnh nghĩa là gì?
\end{enumerate}
\end{checkbox}

% ================= 6. CHMOD =================
\section{\texttt{chmod} --- hai cách đổi quyền}

\subsection{Cách KÝ HIỆU (symbolic) --- dễ đọc cho người}

Cú pháp: \cmd{chmod [ai][+/-/=][quyền] file}

\textbf{Phần ``ai'':}

\begin{center}
\begin{tabular}{|c|>{\raggedright\arraybackslash}p{8cm}|}
\hline
\rowcolor{tablehead}\textbf{Ký tự} & \textbf{Nghĩa} \\
\hline
\cmd{u} & user/owner (chủ) \\
\hline
\cmd{g} & group (nhóm) \\
\hline
\cmd{o} & others (người khác) \\
\hline
\cmd{a} & all (cả ba --- all) \\
\hline
\end{tabular}
\end{center}

\textbf{Phần ``phép toán'':}

\begin{center}
\begin{tabular}{|c|>{\raggedright\arraybackslash}p{8cm}|}
\hline
\rowcolor{tablehead}\textbf{Dấu} & \textbf{Nghĩa} \\
\hline
\cmd{+} & thêm quyền (cộng dồn vào quyền hiện có) \\
\hline
\cmd{-} & bớt quyền \\
\hline
\cmd{=} & \textbf{đặt đúng bằng} (xóa các quyền khác không liệt kê) \\
\hline
\end{tabular}
\end{center}

\begin{notebox}[Phân biệt quan trọng \texttt{+} và \texttt{=}]
\cmd{+x} chỉ \textbf{thêm} quyền \cmd{x}, KHÔNG đụng tới \cmd{r}/\cmd{w} đang có. Còn
\cmd{=} \textbf{đặt tuyệt đối}: \cmd{g=r} nghĩa là group chỉ còn đúng \cmd{r}, mọi
quyền khác của group (kể cả \cmd{w} nếu đang có) bị xóa. Khi muốn \textbf{chắc chắn
khóa} một quyền của ai đó, dùng \cmd{=} (hoặc cách số), đừng dùng \cmd{+}.
\end{notebox}

\textbf{Ví dụ:}
\begin{Verbatim}
chmod u+x file      # them thuc thi cho chu (khong dung quyen khac)
chmod go-w file     # bo quyen ghi cua group VA others
chmod a+r file      # cho tat ca doc duoc
chmod o= file       # xoa sach moi quyen cua others
\end{Verbatim}

\subsection{Cách SỐ (numeric / octal) --- gọn cho dân quen}

\begin{notebox}[Thuật ngữ]
\textit{octal} = hệ \textbf{bát phân}, mỗi chữ số chạy từ \textbf{0 đến 7}. Hệ này
vừa khít với quyền vì mỗi vai chỉ có 3 quyền \cmd{r/w/x} (tối đa \cmd{4+2+1 = 7}).
\end{notebox}

Mỗi vai là một \textbf{chữ số} = tổng của \cmd{r(4) + w(2) + x(1)}:

\begin{center}
\begin{tabular}{|c|c|>{\raggedright\arraybackslash}p{6cm}|}
\hline
\rowcolor{tablehead}\textbf{Số} & \textbf{Quyền} & \textbf{Nghĩa} \\
\hline
7 & \cmd{rwx} & đọc + ghi + chạy (4+2+1) \\
\hline
6 & \cmd{rw-} & đọc + ghi (4+2) \\
\hline
5 & \cmd{r-x} & đọc + chạy (4+1) \\
\hline
4 & \cmd{r-{}-} & chỉ đọc \\
\hline
0 & \cmd{-{}-{}-} & không có gì \\
\hline
\end{tabular}
\end{center}

Một lệnh \cmd{chmod} số gồm \textbf{3 chữ số}: owner, group, others --- theo đúng
thứ tự.

\textbf{Ba con số ``quốc dân'' phải thuộc lòng:}

\begin{center}
\begin{tabular}{|c|c|>{\raggedright\arraybackslash}p{7.5cm}|}
\hline
\rowcolor{tablehead}\textbf{Lệnh} & \textbf{Quyền} & \textbf{Dùng cho} \\
\hline
\cmd{chmod 755} & \cmd{rwxr-xr-x} & \textbf{script \& thư mục cần cho người khác
truy cập}: chủ làm tất, người khác đọc/chạy/vào được \\
\hline
\cmd{chmod 644} & \cmd{rw-r-{}-r-{}-} & \textbf{file thường}: chủ sửa, người khác
chỉ đọc \\
\hline
\cmd{chmod 600} & \cmd{rw-{}-{}-{}-{}-{}-} & \textbf{file riêng tư} (vd khóa SSH):
chỉ mình chủ đọc/ghi \\
\hline
\end{tabular}
\end{center}

\begin{notebox}[Mẹo nhẩm]
755 = ``7 cho mình, 5-5 cho thiên hạ''. 644 = ``6 cho mình, 4-4 cho thiên hạ''.
600 = ``6 cho mình, 0-0 cho thiên hạ'' (khóa kín).
\end{notebox}

\begin{warnbox}[Đừng ``cứ 755 cho mọi thứ'']
755 hợp với script và thư mục \textbf{cần} cho người khác vào. Nhưng với thư
mục/dữ liệu \textbf{riêng tư}, theo tư duy least privilege (mục 10) bạn nên cân nhắc
\cmd{700} (chỉ mình bạn), \cmd{750} (mình + nhóm đọc/vào), hay \cmd{600} cho file
riêng. Cấp đúng mức tối thiểu cần dùng, không cấp dư.
\end{warnbox}

\begin{checkbox}[Tự kiểm tra 5.5]
\begin{enumerate}
\item \cmd{chmod go-w abc} làm gì?
\item Số \cmd{5} tương ứng quyền nào? \cmd{6}? \cmd{7}?
\item \cmd{chmod 600} cho ra chuỗi quyền nào, và phù hợp với loại file gì?
\end{enumerate}
\end{checkbox}

% ================= 7. LAB C =================
\section{Thực hành ngay (lab C): File riêng tư + so sánh hai cách}

Vẫn đang ở \cmd{\textasciitilde/linux-lab/ch5}. (Nếu lỡ thoát, gõ lại
\cmd{cd \textasciitilde/linux-lab/ch5}.)

\begin{labbox}[Lab C --- Bước 1: File riêng tư với \texttt{600}]
\begin{Verbatim}
echo "rieng tu" > bi-mat.txt
chmod 600 bi-mat.txt
ls -l bi-mat.txt
\end{Verbatim}
\textbf{Kết quả mong đợi:}
\begin{Verbatim}
-rw-------@  1 vinh  staff  9 Jun 26 10:05 bi-mat.txt
\end{Verbatim}
Chỉ owner đọc/ghi; group và others \textbf{trống trơn} (\cmd{-{}-{}-{}-{}-{}-}). Đây
đúng là kiểu quyền dùng cho khóa bí mật. (Dấu \cmd{@} nếu có chỉ là thuộc tính mở
rộng của macOS, bỏ qua.)
\end{labbox}

\begin{labbox}[Lab C --- Bước 2: Quan sát \texttt{go-w} có TÁC DỤNG THẬT]
Để thấy rõ \cmd{go-w} \textit{gỡ} quyền ghi, ta cố tình tạo trạng thái có \cmd{w}
cho group và others trước:
\begin{Verbatim}
chmod 664 bi-mat.txt
ls -l bi-mat.txt
\end{Verbatim}
\textbf{Kết quả mong đợi:} group và others đều có \cmd{w}:
\begin{Verbatim}
-rw-rw-r--@  1 vinh  staff  9 Jun 26 10:05 bi-mat.txt
\end{Verbatim}
Giờ gỡ \cmd{w} của group + others:
\begin{Verbatim}
chmod go-w bi-mat.txt
ls -l bi-mat.txt
\end{Verbatim}
\textbf{Kết quả mong đợi:} \cmd{w} của group biến mất:
\begin{Verbatim}
-rw-r--r--@  1 vinh  staff  9 Jun 26 10:05 bi-mat.txt
\end{Verbatim}
\textbf{Quan sát:} lần này bạn \textbf{thấy rõ} \cmd{w} bị gỡ (từ \cmd{rw-rw-} về
\cmd{rw-r-{}-}). Nếu chạy \cmd{go-w} trên một file vốn đã không có \cmd{w} cho
group/others (vd \cmd{644}), lệnh vẫn chạy đúng nhưng \textbf{không có gì thay đổi
nhìn thấy được} --- đừng tưởng lệnh hỏng.

Đặt lại \cmd{600} cho an toàn rồi sang bước sau:
\begin{Verbatim}
chmod 600 bi-mat.txt
\end{Verbatim}
\end{labbox}

\begin{labbox}[Lab C --- Bước 3: So sánh số 644 và 755 trên \texttt{chao.sh}]
\begin{Verbatim}
chmod 644 chao.sh
ls -l chao.sh
\end{Verbatim}
\textbf{Kết quả mong đợi:} mất \cmd{x}, quay về file thường:
\begin{Verbatim}
-rw-r--r--@  1 vinh  staff  26 Jun 26 10:00 chao.sh
\end{Verbatim}

\begin{Verbatim}
chmod 755 chao.sh
ls -l chao.sh
\end{Verbatim}
\textbf{Kết quả mong đợi:} ai cũng chạy/đọc được, chỉ chủ sửa được:
\begin{Verbatim}
-rwxr-xr-x@  1 vinh  staff  26 Jun 26 10:00 chao.sh
\end{Verbatim}

\textbf{Quan sát:} cùng một file, \cmd{644} và \cmd{755} chỉ khác nhau ở các chữ
\cmd{x}. Bạn vừa ``lập trình quyền'' bằng con số. So với \cmd{chmod u+x} ở lab B ---
hai cách \textbf{cho kết quả tương đương}, chỉ là cách viết khác nhau.
\end{labbox}

\begin{labbox}[Lab C --- Bước 4: Đối chiếu lại bằng ký hiệu (tùy chọn)]
\begin{Verbatim}
chmod go-x chao.sh
ls -l chao.sh
\end{Verbatim}
\textbf{Kết quả mong đợi:} bỏ \cmd{x} của group+others, còn lại owner chạy được:
\begin{Verbatim}
-rwxr--r--@  1 vinh  staff  26 Jun 26 10:00 chao.sh
\end{Verbatim}
\end{labbox}

\begin{labbox}[Lab C --- Bước 5: DỌN DẸP (quan trọng)]
\begin{Verbatim}
cd ~
rm -r ~/linux-lab/ch5
\end{Verbatim}
\begin{itemize}
\item \cmd{cd \textasciitilde}: về thư mục nhà (\cmd{\textasciitilde} = home của bạn).
\item \cmd{rm -r}: xóa thư mục cùng mọi thứ bên trong (\cmd{-r} = recursive, đệ quy).
\end{itemize}
\textbf{Kết quả mong đợi:} im lặng = đã xóa sạch sân tập ch5. Kiểm chứng:
\begin{Verbatim}
ls ~/linux-lab
\end{Verbatim}
sẽ không còn thấy \cmd{ch5}.
\end{labbox}

\begin{warnbox}[Trấn an + cảnh báo]
\cmd{rm -r} ở đây chỉ xóa đúng thư mục tập của bạn
(\cmd{\textasciitilde/linux-lab/ch5}), không liên quan gì tới hệ thống --- an toàn.
NHƯNG hãy ghi nhớ nguyên tắc chung: \textbf{tuyệt đối không} gõ \cmd{rm -r} (nhất là
kèm \cmd{sudo}) với đường dẫn bắt đầu bằng \cmd{/} hay \cmd{\textasciitilde} trừ khi
bạn chắc 100\%. Luôn \cmd{ls} đường dẫn để nhìn rõ mình sắp xóa gì \textbf{trước
khi} xóa.
\end{warnbox}

\begin{checkbox}[Tự kiểm tra 5.6]
\begin{enumerate}
\item Sau \cmd{chmod 600 bi-mat.txt}, group và others thấy được gì?
\item Khác biệt hiển thị giữa \cmd{644} và \cmd{755} nằm ở ký tự nào?
\item \cmd{rm -r} khác \cmd{rm} thường ở chỗ nào?
\end{enumerate}
\end{checkbox}

% ================= 8. CHOWN / CHGRP =================
\section{\texttt{chown} / \texttt{chgrp} --- đổi CHỦ và đổi NHÓM (sơ lược)}

Khác với \cmd{chmod} (đổi \textit{quyền}), hai lệnh này đổi \textit{ai sở hữu}:

\begin{itemize}
\item \textbf{\cmd{chown}} = \textit{change owner} --- đổi chủ sở hữu file. \\
  Ví dụ: \cmd{chown vinh tepnao.txt} (đổi chủ thành \cmd{vinh}). \\
  Đổi cả chủ và nhóm: \cmd{chown vinh:staff tepnao.txt}.
\item \textbf{\cmd{chgrp}} = \textit{change group} --- đổi nhóm của file. \\
  Ví dụ: \cmd{chgrp staff tepnao.txt}.
\end{itemize}

\begin{warnbox}[Các lệnh trên chỉ để NHẬN DIỆN --- đừng chạy trong lab]
Đổi chủ sở hữu file sang \textbf{người dùng khác} \textbf{LUÔN} cần quyền root
(\cmd{sudo}) trên cả Linux và macOS --- người dùng thường không thể tùy ý ``cho/nhận''
chủ sở hữu (nếu cho phép thì hết bảo mật). Nếu bạn thử \cmd{chown} sang user khác mà
không có quyền, bạn sẽ gặp lỗi \cmd{Operation not permitted} --- \textbf{đừng} vì thế
mà gõ \cmd{sudo}; bạn không cần làm điều này khi mới học.

Riêng \cmd{chgrp}: người dùng thường có thể đổi nhóm file của mình sang một nhóm
\textbf{mà chính họ là thành viên} mà không cần \cmd{sudo}; đổi sang nhóm khác thì
cần root.
\end{warnbox}

\begin{notebox}[Khác biệt trên macOS]
macOS dùng nhóm \cmd{staff} làm nhóm chính (Linux thường tạo nhóm riêng theo tên
user). Cú pháp \cmd{chown user:group} thì giống nhau. macOS còn có lớp bảo vệ riêng
(vd các thư mục hệ thống bị khóa bởi cơ chế bảo vệ toàn vẹn) nên kể cả \cmd{sudo}
cũng có chỗ không đổi được.
\end{notebox}

\textbf{Phân biệt cốt lõi --- đừng nhầm:}

\begin{center}
\begin{tabular}{|c|>{\raggedright\arraybackslash}p{8cm}|}
\hline
\rowcolor{tablehead}\textbf{Lệnh} & \textbf{Đổi cái gì} \\
\hline
\cmd{chmod} & \textbf{Quyền} (ai được đọc/ghi/chạy) \\
\hline
\cmd{chown} & \textbf{Chủ sở hữu} (file thuộc về ai) \\
\hline
\cmd{chgrp} & \textbf{Nhóm} của file \\
\hline
\end{tabular}
\end{center}

% ================= 9. SUDO =================
\section{\texttt{sudo} --- mượn quyền quản trị trong chốc lát}

\subsection{\texttt{sudo} là gì?}

\textbf{\cmd{sudo}} = \textit{superuser do} = ``làm với tư cách siêu người dùng''. Bạn
đặt \cmd{sudo} trước một lệnh để \textbf{chạy đúng lệnh đó} với quyền quản trị (quyền
của root), thường để:

\begin{itemize}
\item sửa \textbf{file hệ thống} (file của hệ điều hành),
\item \textbf{cài đặt} phần mềm vào khu vực chung của máy,
\item khởi động/dừng dịch vụ hệ thống.
\end{itemize}

Cú pháp:
\begin{Verbatim}
sudo <lenh can quyen cao>
\end{Verbatim}
Lần đầu trong phiên, \cmd{sudo} sẽ \textbf{hỏi mật khẩu}. Trên \textbf{macOS}, đó
chính là \textbf{mật khẩu đăng nhập của CHÍNH bạn} (với điều kiện tài khoản của bạn
thuộc nhóm admin/sudoers). Khi gõ, màn hình \textbf{không hiện ký tự nào} (kể cả dấu
sao); cứ gõ rồi Enter. Sau khi đúng, macOS thường nhớ trong vài phút để bạn không
phải gõ lại liên tục.

\begin{warnbox}[CẢNH BÁO bảo mật --- đọc kỹ]
Chỉ gõ mật khẩu khi \textbf{CHÍNH BẠN} vừa chủ động gõ một lệnh bắt đầu bằng
\cmd{sudo}. Nếu một script lạ, một hướng dẫn trên mạng, hay bất kỳ thứ gì
\textbf{tự dưng} hiện ô hỏi mật khẩu \cmd{sudo}, \textbf{ĐỪNG} gõ --- hãy dừng lại và
kiểm tra xem lệnh đó làm gì. Đây là vector lừa đảo (phishing) phổ biến với người
mới. (Liên hệ với cảnh báo copy-paste ở mục 9.3.)
\end{warnbox}

\begin{notebox}[Trên macOS]
Tài khoản người dùng đầu tiên thường là \textbf{admin} và có quyền dùng \cmd{sudo}.
Tài khoản ``standard'' (chuẩn) thì không. Bạn \textbf{không cần} thử \cmd{sudo} trong
chương này.
\end{notebox}

\subsection{\texttt{sudo} khác \texttt{su} thế nào?}

\begin{itemize}
\item \textbf{\cmd{su}} (\textit{switch user}) = \textbf{chuyển hẳn} sang một người
  dùng khác (mặc định là root) và \textbf{ở lại đó} cho tới khi \cmd{exit}. Giống như
  ``ngồi vào ghế root rồi làm liên tục''.
\item \textbf{\cmd{sudo}} = chỉ \textbf{mượn quyền cho đúng một lệnh} rồi trả lại
  ngay. Giống như ``nhờ sếp ký duyệt một giấy tờ'' thay vì ``chiếm ghế sếp''.
\end{itemize}

\cmd{sudo} an toàn hơn vì \textbf{phạm vi hẹp} (một lệnh) và \textbf{có ghi log} (xem
mục Senior). Trên macOS, việc \textbf{đăng nhập trực tiếp bằng tài khoản root bị vô
hiệu hóa mặc định} (root không có mật khẩu để đăng nhập trực tiếp), NHƯNG root vẫn
tồn tại và \cmd{sudo} vẫn nâng quyền lên root (uid \cmd{0}) để chạy lệnh --- đó chính
là lý do \cmd{sudo} là con đường chuẩn. Nói cách khác: macOS không cho bạn ``đăng
nhập làm root'', nhưng \cmd{sudo} vẫn chạy lệnh \textbf{với tư cách root}.

\subsection{``Quyền lớn, trách nhiệm lớn'' --- cảnh báo cho người mới}

\cmd{sudo} bỏ qua mọi rào quyền. Một lệnh xóa nhầm chạy với \cmd{sudo} có thể
\textbf{phá hệ điều hành}. Nguyên tắc vàng:

\begin{warnbox}[Nguyên tắc vàng khi dùng \texttt{sudo}]
\begin{itemize}
\item \textbf{KHÔNG} gõ \cmd{sudo} chỉ vì thấy lệnh báo lỗi quyền --- hãy hiểu
  \textit{vì sao} nó cần quyền cao trước đã.
\item \textbf{KHÔNG} copy-paste lệnh \cmd{sudo} lạ từ Internet khi chưa hiểu nó làm
  gì.
\item \textbf{Đặc biệt nguy hiểm:} tuyệt đối tránh tổ hợp \cmd{sudo rm -r} với đường
  dẫn bắt đầu bằng \cmd{/} hoặc \cmd{\textasciitilde} (vd \cmd{rm -rf /} có thể xóa
  sạch máy). Luôn \cmd{ls} đường dẫn trước khi xóa.
\item Khi học, gần như \textbf{mọi việc luyện tập đều KHÔNG cần \cmd{sudo}} (như cả
  chương này).
\end{itemize}
\end{warnbox}

\begin{checkbox}[Tự kiểm tra 5.7]
\begin{enumerate}
\item \cmd{sudo} viết tắt của gì? Dùng để làm gì?
\item Trên macOS, \cmd{sudo} hỏi mật khẩu nào? Vì sao gõ không thấy ký tự?
\item \cmd{sudo} khác \cmd{su} ở điểm cốt lõi nào?
\end{enumerate}
\end{checkbox}

% ================= 10. GOC SENIOR =================
\section{Góc Senior (vì sao, production, cạm bẫy)}

\begin{seniorbox}[1. Phân quyền là XƯƠNG SỐNG của bảo mật]
Cả ngành DevOps/Security xoay quanh một câu hỏi: \textit{``Thực thể này có thực sự
cần quyền đó không?''} Nguyên tắc nền là \textbf{least privilege (đặc quyền tối
thiểu)}: cấp \textbf{đúng mức tối thiểu} đủ để làm việc, không hơn. Một tài
khoản/tiến trình bị chiếm cũng chỉ gây hại trong phạm vi quyền nó có.
\end{seniorbox}

\begin{seniorbox}[2. Vì sao \texttt{chmod 777} là NGUY HIỂM]
\cmd{777} = \cmd{rwxrwxrwx} = \textbf{ai cũng đọc/ghi/chạy được}. Nghĩa là bất kỳ
người dùng hay tiến trình nào trên máy đều có thể \textbf{sửa hoặc thay thế} file đó
--- kẻ tấn công có thể thay script của bạn bằng mã độc. Câu ``cứ 777 cho chắc, đỡ lỗi
quyền'' là \textbf{sai lầm kinh điển} che giấu lỗ hổng thật. Cực kỳ tránh
\textbf{\cmd{chmod -R 777}} (đệ quy cả cây thư mục) --- đây là một trong những thói
quen tệ nhất, mở toang toàn bộ dự án.
\end{seniorbox}

\begin{seniorbox}[3. Ví dụ thật: khóa SSH bắt buộc \texttt{600}]
Khóa riêng SSH (private key, dùng để đăng nhập server từ xa) \textbf{buộc} phải có
quyền \cmd{600} (chỉ chủ đọc/ghi). Nếu lỏng hơn (vd group/others đọc được), chương
trình \cmd{ssh} sẽ \textbf{từ chối dùng khóa} và báo lỗi đại loại \textit{``permissions
are too open''}. Đây là ví dụ sống động: phần mềm chủ động \textbf{ép} bạn theo least
privilege.
\end{seniorbox}

\begin{seniorbox}[4. Bộ ba nâng cao (biết tên để sau này tra cứu)]
\begin{itemize}
\item \textbf{setuid / setgid}: cho phép một chương trình chạy với quyền của
  \textit{chủ file} thay vì người gọi (mạnh nhưng nhạy cảm về bảo mật).
\item \textbf{sticky bit}: thường đặt trên thư mục dùng chung (như \cmd{/tmp}) để
  \textbf{chỉ chủ file mới xóa được file của mình}, dù thư mục mở cho nhiều người
  ghi.
\end{itemize}
(Chương này chỉ giới thiệu tên; bạn sẽ học sâu ở chương nâng cao.)
\end{seniorbox}

\begin{seniorbox}[5. \texttt{sudo} có AUDIT (ghi log)]
Mỗi lần \cmd{sudo}, hệ thống ghi lại \textit{ai, lúc nào, chạy lệnh gì}. Trong
production, đây là cơ sở để \textbf{truy vết} khi có sự cố --- và là lý do \cmd{sudo}
được ưa hơn việc ``thành root rồi làm gì cũng được mà không dấu vết''.
\end{seniorbox}

\begin{seniorbox}[6. Liên hệ tương lai (cloud/IAM)]
Mô hình owner/group/others + least privilege chính là \textbf{phiên bản sơ khai} của
\textbf{IAM} (Identity and Access Management) trên cloud (AWS/GCP/Azure): ở đó bạn sẽ
gán ``policy'' cho ``user/role'' theo đúng tư duy \textit{ai được làm gì với tài
nguyên nào}. Học chắc \cmd{chmod}/\cmd{sudo} hôm nay là đặt nền cho IAM ngày mai.
\end{seniorbox}

% ================= 11. HIEU LAM =================
\section{Hiểu lầm thường gặp (đọc kỹ để khỏi sai cả đời)}

\begin{longtable}{|>{\raggedright\arraybackslash}p{5.5cm}|>{\raggedright\arraybackslash}p{8cm}|}
\hline
\rowcolor{tablehead}\textbf{Hiểu lầm} & \textbf{Sự thật} \\
\hline
\endfirsthead
\hline
\rowcolor{tablehead}\textbf{Hiểu lầm} & \textbf{Sự thật} \\
\hline
\endhead
``Cứ \cmd{chmod 777} cho chắc, đỡ lỗi quyền.'' & \textbf{SAI.} \cmd{777} mở toang
cho mọi người ghi $\rightarrow$ lỗ hổng bảo mật. Đặt \textbf{đúng} quyền tối thiểu cần dùng (vd
\cmd{644}, \cmd{755}, \cmd{600}). \\
\hline
``root và sudo là một.'' & \textbf{Không.} \cmd{root} là \textbf{tài khoản} quyền
tối cao. \cmd{sudo} là \textbf{công cụ} để mượn quyền root cho \textbf{một lệnh} rồi
trả lại. Bạn không ``là'' root khi dùng \cmd{sudo}. \\
\hline
``Đổi quyền = đổi chủ sở hữu.'' & \textbf{Khác nhau.} \cmd{chmod} đổi
\textbf{quyền} (đọc/ghi/chạy). \cmd{chown} đổi \textbf{chủ}. Hai việc hoàn toàn
riêng. \\
\hline
``Thiếu \cmd{x} thì cứ sửa nội dung file là chạy được.'' & \textbf{Không.} Nội dung
đúng nhưng thiếu bit \cmd{x} thì hệ thống \textbf{vẫn từ chối chạy} (đã thấy ở lab
B). \\
\hline
``Thư mục có \cmd{r} là vào được rồi.'' & \textbf{Không.} Vào (cd) cần \cmd{x}.
\cmd{r} chỉ cho \textbf{liệt kê tên}. \\
\hline
``\cmd{+x} cũng khóa được quyền ghi của người khác.'' & \textbf{Không.} \cmd{+x} chỉ
\textbf{thêm} \cmd{x}, không đụng \cmd{w}. Muốn khóa quyền ghi của group/others phải
dùng \cmd{=} (đặt tuyệt đối) hoặc cách số. \\
\hline
``macOS không có root.'' & \textbf{Không.} root vẫn tồn tại; chỉ là \textbf{đăng
nhập trực tiếp} bằng root bị tắt mặc định. \cmd{sudo} vẫn chạy lệnh với tư cách
root. \\
\hline
\end{longtable}

% ================= 12. FLASHCARDS =================
\section{Thẻ ghi nhớ (flashcards)}

\begin{notebox}[Cách dùng]
Che cột phải, tự trả lời, rồi mở ra kiểm tra.
\end{notebox}

\begin{longtable}{|>{\raggedright\arraybackslash}p{6cm}|>{\raggedright\arraybackslash}p{7.5cm}|}
\hline
\rowcolor{tablehead}\textbf{Hỏi} & \textbf{Đáp} \\
\hline
\endfirsthead
\hline
\rowcolor{tablehead}\textbf{Hỏi} & \textbf{Đáp} \\
\hline
\endhead
\cmd{whoami} làm gì? & In tên đăng nhập của tôi \\
\hline
\cmd{id} cho biết gì? & uid, gid (nhóm chính của phiên), các nhóm \\
\hline
Ký tự đầu \cmd{d} trong \cmd{ls -l} nghĩa là? & Thư mục (directory) \\
\hline
Dấu \cmd{@} sau cụm quyền trên macOS? & File có thuộc tính mở rộng (extended
attributes), không phải quyền \\
\hline
\cmd{r}, \cmd{w}, \cmd{x} giá trị số? & 4, 2, 1 \\
\hline
\cmd{7} = quyền gì? & \cmd{rwx} (đọc+ghi+chạy) \\
\hline
\cmd{5} = quyền gì? & \cmd{r-x} (đọc+chạy) \\
\hline
\cmd{chmod 644} $\rightarrow$ chuỗi? & \cmd{rw-r-{}-r-{}-} (file thường) \\
\hline
\cmd{chmod 755} $\rightarrow$ chuỗi? & \cmd{rwxr-xr-x} (script \& thư mục cần cho người khác truy
cập; thư mục riêng tư nên xét 700/750) \\
\hline
\cmd{chmod 600} $\rightarrow$ chuỗi? & \cmd{rw-{}-{}-{}-{}-{}-} (file riêng tư/khóa) \\
\hline
\cmd{chmod u+x file} & Thêm quyền chạy cho chủ (không đụng quyền khác) \\
\hline
\cmd{chmod go-w file} & Bỏ quyền ghi của group+others \\
\hline
\cmd{+x} vs \cmd{=}? & \cmd{+x} thêm x (cộng dồn); \cmd{=} đặt tuyệt đối, xóa quyền
không liệt kê \\
\hline
\cmd{x} trên thư mục nghĩa là? & Được \cmd{cd} vào (bước qua cửa) \\
\hline
\cmd{chmod} vs \cmd{chown}? & Đổi quyền vs đổi chủ sở hữu \\
\hline
\cmd{sudo} là gì? & Chạy 1 lệnh với quyền quản trị (superuser do) \\
\hline
\cmd{sudo} vs \cmd{su}? & Mượn quyền 1 lệnh vs chuyển hẳn sang user khác \\
\hline
macOS có root không? & Có; chỉ tắt đăng nhập trực tiếp, sudo vẫn chạy lệnh với tư
cách root \\
\hline
Vì sao khóa SSH cần \cmd{600}? & Lỏng hơn thì \cmd{ssh} từ chối dùng \\
\hline
Vì sao \cmd{777} nguy hiểm? & Ai cũng ghi được $\rightarrow$ lỗ hổng bảo mật \\
\hline
Least privilege nghĩa là? & Cấp đúng quyền tối thiểu đủ dùng \\
\hline
\end{longtable}

% ================= 13. ON NGAT QUANG =================
\section{Kế hoạch ôn ngắt quãng}

\begin{itemize}
\item \textbf{Sau 1 ngày:} Mở Terminal, gõ lại \cmd{whoami}, \cmd{id}, \cmd{groups}.
  Tạo nhanh 1 file rồi \cmd{chmod 600} và \cmd{chmod 644}, đọc lại \cmd{ls -l} đến
  khi giải mã được không cần nhìn bảng (nhớ: dấu \cmd{@} trên macOS chỉ là thuộc tính
  mở rộng, bỏ qua). Đọc lại ``Thẻ ghi nhớ'' một lượt.
\item \textbf{Sau 3 ngày:} Không nhìn sách, tự làm lại \textbf{toàn bộ lab B} (script
  \cmd{chao.sh}: tạo $\rightarrow$ chạy thất bại $\rightarrow$ \cmd{chmod u+x} $\rightarrow$ chạy được). Nhẩm 755/644/600
  từ trí nhớ.
\item \textbf{Sau 1 tuần:} Giải lại tất cả ``Tự kiểm tra'' + ``Bài tập mở rộng''
  không xem đáp án. Tự giải thích thành lời cho người khác (hoặc cho chính mình):
  \textit{``Vì sao chmod 777 nguy hiểm?''} và \textit{``sudo khác su ở đâu?''}. Nếu
  giải thích trôi chảy = đã thuộc.
\end{itemize}

% ================= 14. TOM TAT =================
\section{Tóm tắt 1 trang}

\begin{itemize}
\item \textbf{Đa người dùng $\rightarrow$ cần phân quyền.} Mỗi file có \textbf{owner / group /
  others} + tài khoản tối cao \textbf{root}.
\item \textbf{Tôi là ai:} \cmd{whoami} (tên), \cmd{id} (uid/gid/nhóm), \cmd{groups}.
  macOS: nhóm chính của phiên thường là \cmd{staff}, uid bắt đầu \textasciitilde 501.
\item \textbf{Đọc \cmd{ls -l}} chuỗi 10 ký tự: \cmd{[loại][owner rwx][group
  rwx][others rwx]}. \cmd{-}=file, \cmd{d}=thư mục, \cmd{l}=link (còn \cmd{c/b/p/s}
  ít gặp). Trên macOS có thể thấy thêm \cmd{@} (extended attributes) hoặc \cmd{+}
  (ACL) sau 10 ký tự --- không phải quyền.
\item \textbf{Quyền:} \cmd{r}=4 (đọc), \cmd{w}=2 (ghi), \cmd{x}=1 (chạy file / vào
  thư mục).
\item \textbf{\cmd{chmod} ký hiệu:} \cmd{u/g/o/a} + \cmd{+/-/=} + \cmd{r/w/x}. VD
  \cmd{chmod u+x}, \cmd{chmod go-w}. Nhớ: \cmd{+} cộng dồn, \cmd{=} đặt tuyệt đối
  (xóa quyền không liệt kê).
\item \textbf{\cmd{chmod} số (octal, mỗi vai 0-7):} 3 chữ số = owner/group/others,
  mỗi số là tổng 4+2+1.
  \begin{itemize}
  \item \cmd{755} = \cmd{rwxr-xr-x} (script/thư mục cần truy cập) $\cdot$
    \cmd{644} = \cmd{rw-r-{}-r-{}-} (file thường) $\cdot$
    \cmd{600} = \cmd{rw-{}-{}-{}-{}-{}-} (riêng tư/khóa). Thư mục/file riêng tư cân
    nhắc \cmd{700/750}.
  \end{itemize}
\item \textbf{\cmd{chown}/\cmd{chgrp}} đổi chủ/nhóm --- \textbf{khác} \cmd{chmod} (đổi
  quyền). Đổi chủ sang user khác \textbf{luôn} cần \cmd{sudo}; chỉ để nhận diện,
  không chạy trong lab.
\item \textbf{\cmd{sudo}} = mượn quyền quản trị cho \textbf{một lệnh}; macOS hỏi mật
  khẩu đăng nhập của \textbf{chính bạn} (gõ không hiện ký tự); chỉ gõ khi \textbf{chính
  bạn} chủ động gõ \cmd{sudo}. Khác \cmd{su} (chuyển hẳn user). macOS tắt đăng nhập
  trực tiếp bằng root nhưng \cmd{sudo} vẫn chạy với tư cách root. Dùng dè dặt; tránh
  \cmd{sudo rm -r} với \cmd{/} hoặc \cmd{\textasciitilde}.
\item \textbf{Senior:} least privilege là gốc; \cmd{chmod 777} (nhất là \cmd{-R 777})
  nguy hiểm; khóa SSH bắt buộc \cmd{600}; \cmd{sudo} có log để audit; đây là nền của
  IAM cloud sau này.
\end{itemize}

% ================= 15. BAI TAP =================
\section{Bài tập mở rộng (có đáp án)}

\begin{labbox}[Chuẩn bị sân tập]
Làm trong sân tập an toàn. Trước khi bắt đầu:
\cmd{mkdir -p \textasciitilde/linux-lab/ch5-bt \&\& cd \textasciitilde/linux-lab/ch5-bt}.
Xong nhớ dọn:
\cmd{cd \textasciitilde \&\& rm -r \textasciitilde/linux-lab/ch5-bt}.
\end{labbox}

\textbf{BT1.} Tạo file \cmd{note.txt} rồi đặt quyền sao cho \textbf{chủ đọc+ghi,
group chỉ đọc, others không có gì}. Viết lệnh (cả 2 cách: số và ký hiệu) và cho biết
chuỗi \cmd{ls -l} mong đợi.

\textbf{BT2.} Có file \cmd{run.sh} đang là \cmd{-rw-r-{}-r-{}-}. Hãy làm cho
\textbf{mọi người đều chạy được} nhưng \textbf{chỉ chủ sửa được}. Viết lệnh số, và
lệnh ký hiệu tương đương.

\textbf{BT3.} Giải mã bằng lời: \cmd{drwxr-x-{}-{}-}. Đây là gì, ai làm được gì?

\textbf{BT4.} Vì sao đặt \cmd{chmod 777 \textasciitilde} (cả thư mục nhà) là ý tưởng
tệ? Nêu 2 lý do.

\textbf{BT5.} Bạn gặp lỗi \cmd{ssh} từ chối khóa riêng và than phiền ``permissions too
open''. Lệnh nào sửa, và vì sao?

\subsection*{Đáp án Bài tập mở rộng}

\textbf{BT1.}
\begin{Verbatim}
echo "ghi chu" > note.txt
chmod 640 note.txt          # cach so
# hoac cach ky hieu tuong duong:
chmod u=rw,g=r,o= note.txt
ls -l note.txt
\end{Verbatim}
Chuỗi mong đợi: \cmd{-rw-r-{}-{}-{}-{}-} (owner \cmd{rw-}, group \cmd{r-{}-}, others
\cmd{-{}-{}-} = số \cmd{640}). (Trên macOS có thể có thêm dấu \cmd{@}:
\cmd{-rw-r-{}-{}-{}-{}-@} --- bình thường.)

\textbf{BT2.}
\begin{Verbatim}
chmod 755 run.sh            # cach so -> rwxr-xr-x
# hoac ky hieu, DAT TUYET DOI cho chac chan:
chmod u=rwx,go=rx run.sh
ls -l run.sh
\end{Verbatim}
Kết quả: \cmd{-rwxr-xr-x}. Mọi người chạy/đọc được, chỉ owner có \cmd{w} (sửa).

\begin{notebox}[Lưu ý sư phạm quan trọng]
Đề bài cho file đang là \cmd{-rw-r-{}-r-{}-}. Trong trường hợp đó, \cmd{chmod a+x
run.sh} cũng \textit{ra đúng} \cmd{-rwxr-xr-x} --- nhưng đó chỉ là \textbf{may mắn nhờ
trạng thái ban đầu}. \cmd{+x} chỉ \textbf{thêm} \cmd{x}, KHÔNG đụng tới \cmd{w}; nếu
file ban đầu lỡ có \cmd{w} cho group/others (vd vừa bị \cmd{chmod 666}), thì \cmd{a+x}
sẽ KHÔNG khóa được quyền sửa của người khác $\rightarrow$ vi phạm yêu cầu ``chỉ chủ sửa được''.
Vì vậy, khi cần \textbf{đảm bảo} kết quả, hãy dùng cách số (\cmd{755}) hoặc ký hiệu
\textbf{đặt tuyệt đối} (\cmd{u=rwx,go=rx}), đừng dựa vào \cmd{+x}.
\end{notebox}

\textbf{BT3.} \cmd{drwxr-x-{}-{}-}:
\begin{itemize}
\item Ký tự đầu \cmd{d} $\rightarrow$ \textbf{thư mục}.
\item \cmd{rwx} (owner): chủ đọc nội dung, ghi/tạo/xóa file bên trong, và \cmd{cd}
  vào được.
\item \cmd{r-x} (group): thành viên nhóm đọc danh sách và \cmd{cd} vào được,
  \textbf{không} tạo/xóa file.
\item \cmd{-{}-{}-} (others): người ngoài \textbf{không thấy, không vào} được gì.
\end{itemize}

\textbf{BT4.} \cmd{chmod 777 \textasciitilde} tệ vì:
\begin{enumerate}
\item \textbf{Bảo mật:} mọi người dùng/tiến trình khác đều \textbf{ghi} được vào thư
  mục nhà của bạn $\rightarrow$ có thể thay/cắm file độc, xóa dữ liệu.
\item \textbf{Phá công cụ dựa trên quyền:} nhiều chương trình (vd \cmd{ssh})
  \textbf{từ chối hoạt động} khi thấy thư mục/khóa quyền quá mở, nên thay vì ``đỡ
  lỗi'' bạn lại \textbf{gây thêm lỗi}; đồng thời vi phạm least privilege.
\end{enumerate}

\textbf{BT5.}
\begin{Verbatim}
chmod 600 ~/.ssh/id_ed25519     # (ten khoa co the khac, vd id_rsa)
\end{Verbatim}
Vì khóa riêng SSH \textbf{bắt buộc} chỉ owner đọc/ghi (\cmd{rw-{}-{}-{}-{}-{}-} =
\cmd{600}). Nếu group/others đọc được, \cmd{ssh} coi là không an toàn (người khác có
thể trộm khóa) nên \textbf{từ chối dùng}. Đặt \cmd{600} khôi phục đúng quyền tối
thiểu, \cmd{ssh} chấp nhận lại.

% ================= DAP AN TU KIEM TRA =================
\section*{Đáp án các mục ``Tự kiểm tra''}
\addcontentsline{toc}{section}{Đáp án các mục ``Tự kiểm tra''}

\begin{checkbox}[Đáp án 5.1]
\begin{enumerate}
\item Owner (chủ), group (nhóm), others (người khác).
\item root là superuser --- quyền tối cao, bỏ qua mọi rào quyền, làm được mọi thứ.
\item Vì hệ điều hành vẫn chạy nhiều tài khoản/tiến trình hệ thống cùng lúc; phân
  quyền cô lập chúng để bảo mật và tránh phá nhau.
\end{enumerate}
\end{checkbox}

\begin{checkbox}[Đáp án 5.2]
\begin{enumerate}
\item \cmd{whoami}.
\item \cmd{uid} định danh \textbf{người dùng}; \cmd{gid} định danh \textbf{nhóm
  chính} (primary group) của phiên làm việc.
\item \cmd{staff}.
\end{enumerate}
\end{checkbox}

\begin{checkbox}[Đáp án 5.3]
\begin{enumerate}
\item owner: đọc + ghi (\cmd{rw-}); others: chỉ đọc (\cmd{r-{}-}).
\item \cmd{r=4}, \cmd{w=2}, \cmd{x=1}.
\item Liệt kê được tên file bên trong (nhờ \cmd{r}), nhưng \textbf{không \cmd{cd}
  vào được} và không truy cập nội dung (thiếu \cmd{x}).
\end{enumerate}
\end{checkbox}

\begin{checkbox}[Đáp án 5.4]
\begin{enumerate}
\item Vì file \cmd{chao.sh} chưa có quyền \cmd{x} (đang \cmd{-rw-r-{}-r-{}-}), hệ
  thống không cho chạy.
\item Thêm quyền thực thi (\cmd{x}) cho owner $\rightarrow$ thành \cmd{-rwxr-{}-r-{}-}.
\item \cmd{./} = ``file nằm ngay trong thư mục hiện tại'' (chỉ rõ vị trí để shell
  tìm và chạy).
\end{enumerate}
\end{checkbox}

\begin{checkbox}[Đáp án 5.5]
\begin{enumerate}
\item Bỏ quyền \textbf{ghi (\cmd{w})} của \textbf{group và others} trên file
  \cmd{abc}.
\item \cmd{5}=\cmd{r-x} (đọc+chạy); \cmd{6}=\cmd{rw-} (đọc+ghi); \cmd{7}=\cmd{rwx}
  (đọc+ghi+chạy).
\item \cmd{chmod 600} $\rightarrow$ \cmd{rw-{}-{}-{}-{}-{}-}; phù hợp \textbf{file riêng tư} như
  khóa SSH (chỉ chủ đọc/ghi).
\end{enumerate}
\end{checkbox}

\begin{checkbox}[Đáp án 5.6]
\begin{enumerate}
\item Không thấy gì (group và others là \cmd{-{}-{}-}); chỉ owner đọc/ghi được.
\item Ở các chữ \cmd{x}: \cmd{755} có \cmd{x} cho cả ba vai (\cmd{rwxr-xr-x}),
  \cmd{644} không có \cmd{x} nào (\cmd{rw-r-{}-r-{}-}).
\item \cmd{rm -r} xóa \textbf{đệ quy} cả thư mục và mọi thứ bên trong; \cmd{rm}
  thường chỉ xóa \textbf{một file} (không xóa được thư mục có nội dung).
\end{enumerate}
\end{checkbox}

\begin{checkbox}[Đáp án 5.7]
\begin{enumerate}
\item \textit{superuser do}; dùng để chạy \textbf{một lệnh} với quyền quản trị
  (root), vd sửa file hệ thống/cài đặt.
\item Mật khẩu \textbf{đăng nhập của chính bạn} (với điều kiện tài khoản của bạn
  thuộc nhóm admin) trên macOS; màn hình ẩn hoàn toàn ký tự là \textbf{cố ý} --- để
  người nhìn qua vai (shoulder surfing) không đọc/đoán được mật khẩu; cứ gõ và Enter.
\item \cmd{sudo} mượn quyền cho \textbf{đúng một lệnh} rồi trả lại (hẹp, có log);
  \cmd{su} \textbf{chuyển hẳn} sang user khác và ở lại cho tới khi \cmd{exit}.
\end{enumerate}
\end{checkbox}

\vspace{1cm}
\begin{notebox}[Chuẩn bị cho Chương 6]
Bạn đã kiểm soát được ``ai làm gì'' với file. Chương sau ta sẽ bước vào \textbf{tiến
trình \& quản lý chương trình đang chạy} (xem, tìm, dừng tiến trình) --- nơi khái
niệm ``chủ sở hữu'' và quyền lại tiếp tục phát huy. Hẹn gặp lại!
\end{notebox}

\end{document}